\begin{table*}[t]
\centering
\caption{Component-fault ranking on fault-containing test cycles of the random split (optimistic reference). Top-1 is the share of cycles whose first-ranked component is truly faulty. Static policies are included to test whether the priority rule adds information beyond the severity and effort constants. The leakage-safe comparison is in Tables~\ref{tab:ceilingbound} and~\ref{tab:boundcheck}.}
\label{tab:ranking}
\footnotesize
\begin{tabularx}{\textwidth}{Xrrr}
\toprule
Method & LRAP & Top-1 & MRR \\
\midrule
\multicolumn{4}{l}{\emph{Static policy baselines (do not use classifier probabilities)}} \\
Severity-only policy & 0.781 & 0.697 & 0.821 \\
Low-effort-only policy & 0.805 & 0.735 & 0.845 \\
Severity/effort heuristic & 0.837 & 0.735 & 0.859 \\
\midrule
\multicolumn{4}{l}{\emph{Explanation-only}} \\
Explanation confidence only & 0.984 & 0.986 & 0.993 \\
\midrule
\multicolumn{4}{l}{\emph{Probability-informed policies (all saturate on this split)}} \\
Probability only & 1.000 & 1.000 & 1.000 \\
Permutation weighted & 1.000 & 1.000 & 1.000 \\
Local occlusion & 1.000 & 1.000 & 1.000 \\
Priority score (prob.\ $+$ evidence $+$ severity/effort) & 1.000 & 1.000 & 1.000 \\
\bottomrule
\end{tabularx}
\end{table*}
